\chapter{Velocity Obstacles}
\label{ch:ch12}
% Function names for pseudocode are prefixed with "Vo" to stay unique
% across the book.
\SetKwFunction{VoTimeToCollision}{TimeToCollision}
\SetKwFunction{VoInside}{InsideVO}
\SetKwFunction{VoChoose}{ChooseVelocity}
\SetKwFunction{VoSimulate}{SimulateVO}
\SetKwFunction{VoPreferred}{PreferredVelocity}
\SetKwFunction{VoSamples}{VelocitySamples}

\begin{objectives}
\begin{itemize}
  \item Explain why local collision avoidance is a choice of \emph{velocity} rather than a choice of path, and set up an encounter between two discs in relative coordinates.
  \item Construct the collision cone and the velocity obstacle of a moving disc by hand: apex, tangent rays and the half-angle $\theta=\arcsin\bigl((r_A+r_B)/\norm{\pos_B-\pos_A}\bigr)$.
  \item Compute the time to collision of a candidate velocity from a quadratic equation, and truncate the velocity obstacle with a time horizon $\ttc$.
  \item Choose a velocity outside the union of all velocity obstacles that stays as close as possible to the preferred velocity, by sampling or by projection.
  \item Implement a velocity-obstacle simulator in \python, measure minimum separation and path length, and recognise the oscillation that appears when two agents avoid each other at the same time.
\end{itemize}
\end{objectives}

% ---------------------------------------------------------------------
\section{Why this matters for a drone swarm}
\label{sec:ch12-motivation}

Everything in \cref{ch:ch03,ch:ch04,ch:ch05,ch:ch06,ch:ch07,ch:ch08,ch:ch09,ch:ch10,ch:ch11} plans
\emph{positions}: a path is a sequence of cells, and time enters, if at all, as
a discrete index. Such plans are computed before the flight, and they are only
as good as the map they were computed on. Now picture the moment that the plan
did not foresee. A delivery drone follows its \cbs route across a park at
$1\,\mathrm{m/s}$ when an unknown quadrotor appears $7\,\mathrm{m}$ away,
flying at $0.9\,\mathrm{m/s}$ on a course that crosses the drone's own. The
tracker of \cref{ch:ch18} reports its position and velocity; nothing says that
it will yield. Within five seconds the two aircraft will occupy the same piece
of sky. Replanning the whole route through a graph is the wrong tool at this
time scale: it takes too long, it needs a map that does not yet contain the
intruder, and the intruder will have moved by the time the new route is
ready. What the drone needs is an answer to a much smaller question:
\emph{which velocity should I fly during the next fraction of a second so
that no collision occurs?}

The velocity obstacle (VO) of Fiorini and Shiller~\cite{fiorini1998motion}
answers exactly this question. It looks at the encounter from the drone's own
point of view, in which the intruder is a disc that moves with the
\emph{relative} velocity of the two aircraft. It then describes, as a single
cone in the space of velocities, every velocity of the drone that leads to a
collision if the intruder keeps its current velocity. Avoiding the intruder
means picking a velocity outside the cone, and picking the one that is closest
to the velocity the drone wanted to fly anyway makes the manoeuvre small. The
computation is a handful of dot products per obstacle, which is why VO and its
descendants run at hundreds of hertz on board of small vehicles.

In the four-layer architecture of \cref{ch:ch24}, velocity obstacles are the
\emph{local safety layer}. The global planner (\cbs or \ecbs,
\cref{ch:ch09,ch:ch10}) produces nominal routes; the prediction layer
(\cref{ch:ch18,ch:ch20}) delivers, for every tracked object, an estimated
position and velocity; the local layer takes over the velocity command whenever
a predicted collision is closer than a safety horizon and hands control back
once the danger has passed. The reciprocal variants of \cref{ch:ch13} (RVO and
\orca) are what the swarm members use among themselves; the plain velocity
obstacle of this chapter is what a drone uses against an object that does not
cooperate.

The chapter proceeds as follows. \Cref{sec:ch12-intuition} gives the idea in
one picture. \Cref{sec:ch12-problem} derives the geometry: relative motion,
the collision cone and its tangent half-angle, the velocity obstacle, the time
to collision and the truncated velocity obstacle. \Cref{sec:ch12-algorithm}
turns the geometry into three short algorithms (membership test, velocity
selection, simulation), \cref{sec:ch12-example} runs them on the crossing
encounter above with every number produced by the chapter's code, and
\cref{sec:ch12-properties} states what velocity obstacles guarantee and what
they do not. \Cref{sec:ch12-oscillation} shows the oscillation that arises when
two agents avoid each other simultaneously, together with an experiment on
crossing agents; it is the motivation for \cref{ch:ch13}.
\Cref{sec:ch12-variants} covers multiple and static obstacles and
\cref{sec:ch12-3d} the cone in three dimensions. \Cref{sec:ch12-implementation}
and \cref{sec:ch12-drone} discuss the implementation and the place of VO in the
drone system.

% ---------------------------------------------------------------------
\section{The idea in plain words}
\label{sec:ch12-intuition}

Two discs move through the plane with constant velocities. Whether they will
collide does not depend on where either of them is going in absolute terms; it
depends only on how one moves \emph{relative} to the other. So sit on drone
$A$. From there, drone $B$ is an object at the relative position
$\pos_B - \pos_A$ that drifts with the relative velocity $\vel_B - \vel_A$.
Two more simplifications make the picture crisp. First, by the Minkowski-sum
argument of \cref{ch:ch02}, two discs of radii $r_A$ and $r_B$ overlap exactly
when their centres are closer than $r_A + r_B$; so we may shrink $A$ to a point
and fatten $B$ to a disc of radius $r_A + r_B$. Second, instead of letting the
fat disc drift towards the point, let the point move towards the fixed disc
with the velocity $\vel_A - \vel_B$; the two descriptions are the same motion
seen from different frames.

The question ``will they collide?'' has now become a question about a ray: the
point starts at the origin and travels along the ray in the direction of the
relative velocity. It hits the disc if and only if the ray intersects the disc
(\cref{fig:ch12-idea}). The directions whose rays hit the disc form a wedge
between the two tangent lines from the origin to the disc. This wedge is the
\textbf{collision cone}\index{collision cone}. A relative velocity inside it
leads to a collision, sooner if the velocity is large and later if it is small;
a relative velocity outside it never does, however fast.

\begin{figure}[tb]
  \centering
  \inputfigure{ch12/idea}
  \caption{The encounter seen from $A$. Agent $A$ is a point at the origin,
  agent $B$ is the disc of radius $r_A+r_B$ (its own radius is drawn inside)
  at the relative position $\pos_B-\pos_A$. A relative velocity
  $\vel_A-\vel_B$ is dangerous exactly when its ray from the origin meets the
  disc: the red ray hits, the grey ray grazes the disc along a tangent, the
  green ray misses. The shaded wedge between the tangents is the collision
  cone.}
  \label{fig:ch12-idea}
\end{figure}

The final step translates the picture back to the velocities that $A$ can
actually command. Since the relative velocity is $\vel_A - \vel_B$, the
condition ``$\vel_A - \vel_B$ lies in the cone'' is the condition ``$\vel_A$
lies in the cone shifted by $\vel_B$''. The shifted cone is the
\textbf{velocity obstacle}\index{velocity obstacle}: a wedge in $A$'s velocity
space with its apex at $\vel_B$. To avoid $B$, agent $A$ chooses any velocity
outside the wedge, and among those it chooses the one closest to the velocity
it prefers. Two refinements make this practical. A collision that lies far in
the future need not be avoided yet, so the wedge is cut off by a time horizon.
And with several obstacles, $A$ avoids the union of their wedges.

\begin{keyidea}
Look at an encounter from the agent's own frame: the other agent is a disc of
radius $r_A + r_B$, and the relative velocity $\vel_A - \vel_B$ is dangerous
exactly when its ray hits that disc. The dangerous relative velocities form a
cone with half-angle $\arcsin\bigl((r_A+r_B)/\norm{\pos_B-\pos_A}\bigr)$; the
velocity obstacle is this cone translated to the apex $\vel_B$. Avoiding a
collision means choosing a velocity outside the velocity obstacle, as close as
possible to the preferred one.
\end{keyidea}

% ---------------------------------------------------------------------
\section{Problem statement and notation}
\label{sec:ch12-problem}

\subsection{Relative motion}

Throughout, an \textbf{agent}\index{agent!disc model} is a disc that moves as
a point mass with a commanded velocity (the single-integrator model of
\cref{ch:ch02}): its centre obeys $\pos(t) = \pos + t\,\vel$ as long as the
velocity $\vel$ is held. Agent $A$ has centre $\pos_A$, velocity $\vel_A$ and
radius $r_A$; agent $B$ (a drone, a bird, a static pole) has $\pos_B$, $\vel_B$,
$r_B$. Velocities may be changed instantaneously; acceleration limits are
discussed in \cref{sec:ch12-variants} and in \cref{ch:ch14}.

\begin{definition}[Relative position and velocity]
\label{def:ch12-relative}
The \textbf{relative position}\index{relative position} of $B$ with respect to
$A$ and the \textbf{relative velocity}\index{relative velocity} of $A$ with
respect to $B$ are
\begin{equation}
  \pos_{\mathrm{rel}} = \pos_B - \pos_A, \qquad
  \vel_{\mathrm{rel}} = \vel_A - \vel_B .
  \label{eq:ch12-relative}
\end{equation}
The \textbf{combined radius}\index{combined radius} is $R = r_A + r_B$ and the
current distance is $d = \norm{\pos_{\mathrm{rel}}}$. We assume $d > R$: the
agents are apart now.
\end{definition}

The sign conventions matter and are chosen so that $A$ is the mover. If both
agents keep their velocities, the centre of $B$ seen from the centre of $A$ is
\begin{equation}
  \pos_B(t) - \pos_A(t)
  = (\pos_B + t\vel_B) - (\pos_A + t\vel_A)
  = \pos_{\mathrm{rel}} - t\,\vel_{\mathrm{rel}} ,
  \label{eq:ch12-relmotion}
\end{equation}
which is the position of a \emph{fixed} disc centre $\pos_{\mathrm{rel}}$ seen
from a point that has moved to $t\,\vel_{\mathrm{rel}}$. By the Minkowski-sum
proposition for two discs in \cref{ch:ch02}, the discs overlap at time $t$ if
and only if $\norm{\pos_B(t)-\pos_A(t)} \le R$, that is,
\begin{equation}
  \text{$A$ and $B$ collide at time $t\ge0$}
  \iff \norm{\pos_{\mathrm{rel}} - t\,\vel_{\mathrm{rel}}} \le R
  \iff t\,\vel_{\mathrm{rel}} \in \disc{\pos_{\mathrm{rel}}}{R},
  \label{eq:ch12-collision}
\end{equation}
where $\disc{\vect{c}}{r}$ is the closed disc of radius $r$ around $\vect{c}$.
The disc $\disc{\pos_{\mathrm{rel}}}{R} = \disc{\pos_{\mathrm{rel}}}{r_B}
\oplus \disc{\vect{0}}{r_A}$ is the Minkowski sum\index{Minkowski sum!in
velocity obstacles} of $B$'s disc with $A$'s disc placed at the origin: $A$ has
become a point, $B$ has grown by $r_A$.

\subsection{The collision cone}

\begin{definition}[Collision cone]
\label{def:ch12-collision-cone}
The \textbf{collision cone}\index{collision cone|textbf} of $A$ induced by $B$ is
the set of relative velocities that lead to a collision at some future time,
\begin{equation}
  CC_{A|B} = \set{\vel_{\mathrm{rel}} \in \R^2 :
    \exists\, t > 0 \ \text{with}\
    t\,\vel_{\mathrm{rel}} \in \disc{\pos_{\mathrm{rel}}}{R}}.
  \label{eq:ch12-cc}
\end{equation}
Equivalently, $\vel_{\mathrm{rel}} \in CC_{A|B}$ if and only if the ray
$\set{t\,\vel_{\mathrm{rel}} : t > 0}$ from the origin intersects the disc
$\disc{\pos_{\mathrm{rel}}}{R}$.
\end{definition}

The cone depends on $\pos_{\mathrm{rel}}$ and $R$ only, not on any velocity.
It contains, with any vector, every positive multiple of that vector: if the
ray in direction $\vel$ hits the disc, so does the ray in direction $2\vel$,
only sooner. This is what makes $CC_{A|B}$ a cone, and it is why the whole set
is described by an angle.

\begin{theorem}[Tangent half-angle of the collision cone]
\label{thm:ch12-cone}
Let $d = \norm{\pos_{\mathrm{rel}}} > R$ and let
\begin{equation}
  \theta = \arcsin\frac{R}{d} \in \Bigl(0, \frac{\pi}{2}\Bigr).
  \label{eq:ch12-halfangle}
\end{equation}
Then $CC_{A|B}$ is the closed cone of all non-zero vectors that make an angle
of at most $\theta$ with $\pos_{\mathrm{rel}}$:
$CC_{A|B} = \set{\vel \ne \vect{0} : \angle(\vel, \pos_{\mathrm{rel}}) \le \theta}$.
Its two boundary rays are the tangent rays from the origin to the circle of
radius $R$ around $\pos_{\mathrm{rel}}$; they touch the circle at distance
$\sqrt{d^2 - R^2}$ from the origin, and their directions are
$\pos_{\mathrm{rel}}/d$ rotated by $+\theta$ and by $-\theta$.
\end{theorem}
\begin{proof}
Let $\vel \ne \vect{0}$, write $\vect{u} = \vel/\norm{\vel}$ and let $\varphi
\in [0,\pi]$ be the angle between $\vect{u}$ and $\pos_{\mathrm{rel}}$. The ray
$\set{t\vect{u} : t \ge 0}$ meets the disc if and only if the distance from the
centre $\pos_{\mathrm{rel}}$ to the ray is at most $R$. If $\varphi \le \pi/2$,
the point of the line $\set{t\vect{u} : t \in \R}$ closest to the centre is at
$t^* = \pos_{\mathrm{rel}}\cdot\vect{u} = d\cos\varphi \ge 0$, hence on the ray,
and its distance to the centre is $d\sin\varphi$ (the perpendicular from the
centre to the line; \cref{ch:ch02}, distance from a point to a segment). If
$\varphi > \pi/2$, the closest point of the ray is the origin itself, at
distance $d > R$, and the ray misses. Therefore the ray meets the disc if and
only if $\varphi \le \pi/2$ and $d\sin\varphi \le R$, i.e.\ if and only if
$\varphi \le \arcsin(R/d) = \theta$, because $\sin$ is increasing on
$[0,\pi/2]$ and $R/d < 1$. This proves the description of the cone. A ray with
$\varphi = \theta$ has distance exactly $R$ to the centre, so it touches the
circle at one point $\vect{t}$; the radius to $\vect{t}$ is perpendicular to the
ray, and the right triangle with vertices $\vect{0}$, $\vect{t}$,
$\pos_{\mathrm{rel}}$ has hypotenuse $d$ and legs $R$ and
$\norm{\vect{t}} = \sqrt{d^2-R^2}$, with $\sin\theta = R/d$ at the origin. These
are the tangent points of \cref{ch:ch02}. \qedhere
\end{proof}

The theorem is the whole geometry of this chapter, so it is worth reading
\cref{eq:ch12-halfangle} twice. The cone is narrow when the obstacle is far
($R/d$ small) and opens as the obstacle approaches; when $d$ shrinks to $R$ the
half-angle reaches $90^\circ$ and every velocity that has any component
towards $B$ is dangerous. Forgetting to add $r_A$ to $r_B$ is the most common
mistake in this construction; it makes the cone too narrow and the agent
grazes the obstacle (\cref{sec:ch12-implementation}).

\subsection{The velocity obstacle}

\begin{definition}[Velocity obstacle]
\label{def:ch12-vo}
The \textbf{velocity obstacle}\index{velocity obstacle|textbf} of $A$ induced
by $B$ moving with velocity $\vel_B$ is the collision cone translated by
$\vel_B$,
\begin{equation}
  \VO_{A|B}(\vel_B) = \vel_B \oplus CC_{A|B}
  = \set{\vel_B + \vel : \vel \in CC_{A|B}}
  = \set{\vel_A \in \R^2 : \vel_A - \vel_B \in CC_{A|B}} .
  \label{eq:ch12-vo}
\end{equation}
A velocity $\vel_A$ is \textbf{dangerous} if $\vel_A \in \VO_{A|B}(\vel_B)$
and \textbf{safe} otherwise. When $\vel_B$ is clear from the context we write
$\VO_{A|B}$.
\end{definition}

The velocity obstacle lives in the velocity space of $A$: its points are
velocities that $A$ could command, and $\VO_{A|B}$ is the set of those that
$A$ must not command. It is a wedge with its \textbf{apex}\index{velocity
obstacle!apex} at $\vel_B$, its axis parallel to $\pos_{\mathrm{rel}}$ and the
half-angle $\theta$ of \cref{thm:ch12-cone}. Unlike the collision cone, it is
not closed under scaling unless $\vel_B = \vect{0}$, and it changes whenever
$B$ changes its velocity. \Cref{fig:ch12-geometry} shows both pictures for the
encounter of the introduction, which is also the worked example of
\cref{sec:ch12-example}: in the workspace on the left, the cone of dangerous
relative directions sits at $A$ and the inflated disc around $B$; in the
velocity space on the right, the same wedge sits at the apex $\vel_B$, and the
preferred velocity $\vel_A^{\mathrm{pref}}$ of $A$ falls inside it.

\begin{figure}[tb]
  \centering
  \inputfigure{ch12/geometry}
  \caption{The velocity obstacle of the worked example. (a)~Workspace: $A$ at
  $(-5,0)$ with $r_A=0.5$ flies east at $1\,\mathrm{m/s}$, $B$ at $(0,-5)$
  with $r_B=0.5$ flies north at $0.9\,\mathrm{m/s}$. The dashed circle is the
  inflated disc of radius $R=1$; the tangents from the centre of $A$ enclose
  the collision cone of half-angle $\theta = \arcsin(1/\sqrt{50}) = 8.13^\circ$
  around the direction to $B$. (b)~Velocity space of $A$: the same wedge with
  its apex at $\vel_B=(0,0.9)$ is $\VO_{A|B}$; its legs have the directions
  $(0.8,-0.6)$ and $(0.6,-0.8)$. The preferred velocity
  $\vel_A^{\mathrm{pref}}=(1,0)$ lies inside, so $A$ must deviate. The small
  circle is the truncation disc for $\ttc=5\,\mathrm{s}$
  (\cref{sec:ch12-truncated}); only the darker region is forbidden. The
  velocity $\vel^*$ is the point of the boundary closest to
  $\vel_A^{\mathrm{pref}}$ and $\vel_A'$ is the sampled choice of
  \cref{sec:ch12-example}. The dotted circle bounds the reachable velocities
  $\norm{\vel}\le v_{\max}=1.5$.}
  \label{fig:ch12-geometry}
\end{figure}

\subsection{Time to collision}
\label{sec:ch12-ttc}

Membership in the cone says \emph{whether} a velocity collides. For the
truncation of the next section, and for choosing between bad options, we need
to know \emph{when}.

\begin{definition}[Time to collision]
\label{def:ch12-ttc}
The \textbf{time to collision}\index{time to collision|textbf} of a relative
velocity $\vel_{\mathrm{rel}}$ is the first time at which the discs touch,
\begin{equation}
  t_c(\vel_{\mathrm{rel}}) = \min\set{t \ge 0 :
    \norm{\pos_{\mathrm{rel}} - t\,\vel_{\mathrm{rel}}} \le R},
  \label{eq:ch12-ttc-def}
\end{equation}
with $t_c = \infty$ if the set is empty. For an absolute velocity $\vel_A$ of
$A$ we write $t_c(\vel_A - \vel_B)$.
\end{definition}

Squaring the condition $\norm{\pos_{\mathrm{rel}} - t\,\vel_{\mathrm{rel}}} = R$
gives a quadratic equation in $t$. With the abbreviations $\pos = \pos_{\mathrm{rel}}$ and $\vel = \vel_{\mathrm{rel}}$,
\begin{equation}
  (\vel\cdot\vel)\,t^2 - 2\,(\pos\cdot\vel)\,t + (\pos\cdot\pos - R^2) = 0,
  \qquad
  \frac{\Delta}{4} = (\pos\cdot\vel)^2 - (\vel\cdot\vel)(\pos\cdot\pos - R^2).
  \label{eq:ch12-quadratic}
\end{equation}
This is the ray--circle intersection of \cref{ch:ch02} with the ray starting at
the origin in the direction $\vel$ and the circle centred at $\pos$.

\begin{proposition}[Time to collision]
\label[proposition]{thm:ch12-ttc}
Assume $d > R$ and $\vel \ne \vect{0}$. Then $\vel \in CC_{A|B}$ if and only if
$\Delta \ge 0$ and $\pos\cdot\vel > 0$, and in that case
\begin{equation}
  t_c(\vel) = \frac{\pos\cdot\vel - \sqrt{\Delta/4}}{\vel\cdot\vel} > 0,
  \label{eq:ch12-ttc}
\end{equation}
the smaller root of \cref{eq:ch12-quadratic}. The discs overlap exactly during
the interval between the two roots. If $\vel = \vect{0}$ there is no collision.
\end{proposition}
\begin{proof}
The function $t \mapsto \norm{\pos - t\vel}^2 - R^2$ is a parabola opening
upwards with the value $d^2 - R^2 > 0$ at $t = 0$. It is non-positive exactly
between its two real roots, if they exist, so the discs overlap exactly on that
interval, and a future collision exists if and only if the roots are real
($\Delta \ge 0$) and the smaller one is positive. The product of the roots is
$(\pos\cdot\pos - R^2)/(\vel\cdot\vel) > 0$, so the roots have the same sign,
which is the sign of their sum $2(\pos\cdot\vel)/(\vel\cdot\vel)$. Hence both
are positive if and only if $\pos\cdot\vel > 0$, and then the first contact is
at the smaller root \cref{eq:ch12-ttc}. For $\vel = \vect{0}$ the distance stays
$d > R$. \qedhere
\end{proof}

The discriminant condition is the cone condition in disguise: with
$\varphi$ the angle between $\vel$ and $\pos$, a short computation gives
$\Delta/4 = \norm{\vel}^2\bigl(R^2 - d^2\sin^2\varphi\bigr)$, so
$\Delta \ge 0$ is exactly $d\sin\varphi \le R$ as in the proof of
\cref{thm:ch12-cone}, and $\pos\cdot\vel > 0$ excludes the mirror-image
directions that point away from the disc. A velocity with $\Delta = 0$ is
grazing: the discs touch at a single instant. For the encounter of
\cref{fig:ch12-geometry}, $\pos = (5,-5)$, $R = 1$ and the preferred velocity
gives $\vel = (1,0)-(0,0.9) = (1,-0.9)$; then
$\vel\cdot\vel = 1.81$, $\pos\cdot\vel = 9.5$, $\pos\cdot\pos - R^2 = 49$,
$\Delta/4 = 90.25 - 88.69 = 1.56$ and
$t_c = (9.5 - 1.249)/1.81 = 4.559\,\mathrm{s}$; the discs would overlap from
$t = 4.56$ to $t = 5.94\,\mathrm{s}$.

\subsection{The truncated velocity obstacle}
\label{sec:ch12-truncated}

Taken literally, \cref{def:ch12-vo} forbids a velocity that would cause a
collision in ten minutes with a drone at the other end of the city. Fiorini
and Shiller already distinguished \emph{imminent} collisions from those that
lie beyond a time horizon~\cite{fiorini1998motion}, and every practical
implementation cuts the cone off.

\begin{definition}[Truncated velocity obstacle]
\label{def:ch12-truncated}
For a \textbf{time horizon}\index{time horizon!velocity obstacle} $\ttc > 0$,
the \textbf{truncated velocity obstacle}\index{velocity obstacle!truncated} is
the set of velocities that collide within the horizon,
\begin{equation}
  \VO^{\ttc}_{A|B}(\vel_B) = \set{\vel_A : t_c(\vel_A - \vel_B) \le \ttc}
  \subseteq \VO_{A|B}(\vel_B) .
  \label{eq:ch12-truncated}
\end{equation}
\end{definition}

\begin{proposition}[Shape of the truncated velocity obstacle]
\label[proposition]{thm:ch12-truncated}
In relative-velocity coordinates the truncated cone is a union of discs,
\begin{equation}
  \VO^{\ttc}_{A|B}(\vel_B) - \vel_B
  = \bigcup_{0 < t \le \ttc} \disc{\pos_{\mathrm{rel}}/t}{R/t} .
  \label{eq:ch12-union}
\end{equation}
Every disc $\disc{\pos_{\mathrm{rel}}/t}{R/t}$ is inscribed in the collision
cone (tangent to both legs), and the discs shrink and move towards the apex as
$t$ grows. Consequently $\VO^{\ttc}_{A|B}$ is the cone with its tip cut off by
the disc $\disc{\vel_B + \pos_{\mathrm{rel}}/\ttc}{R/\ttc}$: its boundary consists
of the two legs beyond their points of tangency with this disc and of the arc
of the disc that faces the apex.
\end{proposition}
\begin{proof}
By \cref{eq:ch12-collision}, $\vel$ collides at time $t$ if and only if
$\norm{\pos_{\mathrm{rel}} - t\vel} \le R$, i.e.\ if and only if
$\norm{\pos_{\mathrm{rel}}/t - \vel} \le R/t$, i.e.\ $\vel \in
\disc{\pos_{\mathrm{rel}}/t}{R/t}$. Since $t_c(\vel) \le \ttc$ means that
$\vel$ collides at some $t \le \ttc$, \cref{eq:ch12-union} follows. The disc for
$t$ is the image of $\disc{\pos_{\mathrm{rel}}}{R}$ under the scaling
$\vect{x} \mapsto \vect{x}/t$ about the origin. A scaling about the origin maps
every ray from the origin onto itself and preserves tangency, so the two
tangent rays of $\disc{\pos_{\mathrm{rel}}}{R}$ are also the tangent rays of
every scaled disc. Along a fixed ray inside the cone, the times at which the
point $t\vel$ is inside the original disc form one interval $[t_c, t_2]$, so
the union over $t \le \ttc$ contains the ray's points beyond the first entry
into the disc for $t = \ttc$ and nothing before it; this describes the
boundary. \qedhere
\end{proof}

\begin{figure}[tb]
  \centering
  \inputfigure{ch12/truncated}
  \caption{Truncation, drawn in relative-velocity coordinates for
  $\pos_{\mathrm{rel}}=(4,0)$ and $R=1$ (half-angle $14.5^\circ$). The relative
  velocity $\vel$ collides at time $t$ exactly when $\vel$ lies in the disc
  $\disc{\pos_{\mathrm{rel}}/t}{R/t}$; the discs for $t=1$, $2$ and $4$ are
  drawn, each inscribed in the cone. With the horizon $\ttc=2$ only the dark
  region beyond the disc for $t=2$ is forbidden: the velocity $\vel$, which
  collides at $t_c=3$, is allowed, $\vel'$ with $t_c=1.5$ is not.}
  \label{fig:ch12-truncated}
\end{figure}

\Cref{fig:ch12-truncated} shows the construction and \cref{fig:ch12-geometry}(b)
its instance in the worked example, where the truncation disc for
$\ttc = 5\,\mathrm{s}$ has centre $\vel_B + \pos_{\mathrm{rel}}/5 = (1,-0.1)$
and radius $R/5 = 0.2$. The horizon has a clean interpretation: $A$ ignores
collisions that are more than $\ttc$ seconds away, on the grounds that either
agent will have changed its velocity before then and a fresh decision will be
made at the next time step. It also has a cost: a velocity just outside
$\VO^{\ttc}$ collides in $\ttc + \epsilon$ seconds, so the horizon must be long
enough for $A$ to make the manoeuvre that the next decision will demand.
Practical values are a few seconds for slow indoor vehicles and tens of seconds
for fast aircraft; \cref{sec:ch12-implementation} returns to the choice.

\subsection{Choosing a velocity}
\label{sec:ch12-choosing}

The velocity obstacle tells $A$ what not to do. The decision of what to do
instead needs two more ingredients: what $A$ can fly, and what $A$ wants.

\begin{definition}[Reachable, preferred and feasible velocities]
\label{def:ch12-feasible}
The \textbf{reachable velocities}\index{reachable velocities} of $A$ are the
disc $\mathcal{V}_A = \disc{\vect{0}}{v_{\max}}$ of speeds up to the limit
$v_{\max}$ of \cref{ch:ch02}. The \textbf{preferred
velocity}\index{preferred velocity} $\vel_A^{\mathrm{pref}}$ is the velocity
$A$ would fly if nothing were in the way; in this chapter it points from
$\pos_A$ towards the goal at the nominal speed. Given the obstacles
$B_1,\dots,B_k$ with velocities $\vel_{B_j}$, the \textbf{feasible
velocities}\index{feasible velocities} are
\begin{equation}
  \mathcal{F}_A = \mathcal{V}_A \setminus
    \bigcup_{j=1}^{k} \VO^{\ttc}_{A|B_j}(\vel_{B_j}),
  \label{eq:ch12-feasible}
\end{equation}
and $A$ selects the feasible velocity closest to its preferred one,
\begin{equation}
  \vel_A^{*} = \argmin_{\vel \in \mathcal{F}_A}\ \norm{\vel - \vel_A^{\mathrm{pref}}} .
  \label{eq:ch12-cost}
\end{equation}
\end{definition}

\begin{figure}[tb]
  \centering
  \inputfigure{ch12/feasible}
  \caption{Choosing a velocity with two obstacles. The blue disc is the
  reachable set $\norm{\vel}\le v_{\max}$; the dots are the sampled candidate
  velocities of \cref{alg:ch12-choose}. The orange wedge with its apex at
  $\vel_{B_1}$ is the truncated velocity obstacle of a moving agent, the wedge
  with its apex at the origin that of a static obstacle ($\vel_{B_2}=\vect{0}$).
  The feasible set is what remains of the disc. The preferred velocity
  $\vel^{\mathrm{pref}}$ is forbidden; the chosen velocity $\vel^*$ is the
  feasible point closest to it, here on the truncation arc of the first
  obstacle.}
  \label{fig:ch12-feasible}
\end{figure}

\Cref{fig:ch12-feasible} shows the sets. Three remarks complete the
definition. First, the union in \cref{eq:ch12-feasible} is how
\textbf{multiple obstacles}\index{velocity obstacle!union of} are handled: a
velocity is safe only if it is safe with respect to every obstacle, so the
forbidden set is the union of the individual wedges, and the feasible set can
be small and disconnected. Second, a \textbf{static
obstacle}\index{static obstacle!as velocity obstacle} is the special case
$\vel_B = \vect{0}$: its velocity obstacle is the collision cone itself, with
the apex at the origin, and a polygonal obstacle is handled by the cone
spanned by its extreme tangents (\cref{sec:ch12-variants}). Third, the cost in
\cref{eq:ch12-cost} is the simplest sensible one: it prefers small deviations
in speed and heading equally. Other costs are possible and easy to use with the
sampling algorithm below, for example one that penalises changes from the
\emph{current} velocity to reduce chattering, or one that maximises the speed
towards the goal, which is the heuristic of the original
paper~\cite{fiorini1998motion}.

Two things can go wrong with \cref{eq:ch12-cost}. The feasible set may be
empty, because a fast obstacle is so close that every reachable velocity
collides within $\ttc$; then there is no safe choice and $A$ must pick the
least dangerous velocity, for example the one that maximises $t_c$ or the one
that minimises $\norm{\vel - \vel^{\mathrm{pref}}} + \kappa/t_c(\vel)$ for a
weight $\kappa > 0$. And the minimiser may be hard to compute exactly when
several wedges overlap; the practical answer is to evaluate the cost on a
finite set of candidate velocities, as the algorithm of the next section does,
or to project onto one wedge at a time as \orca does in \cref{ch:ch13}.

% ---------------------------------------------------------------------
\section{The algorithm}
\label{sec:ch12-algorithm}

Velocity obstacles are not a planner with a start and a goal; they are a rule
that is applied at every control step. Three pieces make up the rule: the
membership test of a velocity (\cref{alg:ch12-ttc}), the selection of a
velocity for one agent (\cref{alg:ch12-choose}), and the loop that applies the
selection to every agent and advances the simulation (\cref{alg:ch12-simulate}).

\begin{algorithm}[htb]
\caption{Time to collision and membership in the truncated velocity obstacle}
\label{alg:ch12-ttc}
\KwIn{relative position $\pos$, relative velocity $\vel$, combined radius $R$, horizon $\ttc$}
\KwOut{$t_c(\vel)$, and whether $\vel_A \in \VO^{\ttc}_{A|B}$}
\Fn{\VoTimeToCollision{$\pos$, $\vel$, $R$}}{
  $a \gets \vel\cdot\vel$; $b \gets \pos\cdot\vel$; $c \gets \pos\cdot\pos - R^2$\;
  \If{$c \le 0$}{\KwRet{$0$}\tcp*{already overlapping}}
  \If{$a = 0$}{\KwRet{$\infty$}\tcp*{no relative motion}}
  $\Delta' \gets b^2 - a\,c$\label{alg:ch12-ttc:disc}\;
  \If{$\Delta' < 0$}{\KwRet{$\infty$}\tcp*{the ray misses the disc}}
  $t \gets (b - \sqrt{\Delta'})/a$\label{alg:ch12-ttc:root}\;
  \lIf{$t < 0$}{\KwRet{$\infty$}\tcp*[f]{the disc lies behind $A$}}
  \KwRet{$t$}\;
}
\Fn{\VoInside{$\pos_A$, $r_A$, $\pos_B$, $\vel_B$, $r_B$, $\vel_A$, $\ttc$}}{
  $t \gets$ \VoTimeToCollision{$\pos_B-\pos_A$, $\vel_A-\vel_B$, $r_A+r_B$}\label{alg:ch12-ttc:rel}\;
  \KwRet{$t \le \ttc$}\;
}
\end{algorithm}

\paragraph{Walkthrough of \cref{alg:ch12-ttc}.}
The function \VoTimeToCollision evaluates \cref{eq:ch12-ttc} with the guards
that a real implementation needs. The constant term $c$ is negative when the
discs already overlap; then the collision is happening now and the function
returns $0$, leaving it to the caller to decide what to do (the code of
\cref{sec:ch12-implementation} forbids every velocity that approaches the
other agent). Line~\ref{alg:ch12-ttc:disc} computes a quarter of the
discriminant of \cref{eq:ch12-quadratic}; a negative value means the ray misses
the disc and the velocity is safe for ever. Line~\ref{alg:ch12-ttc:root} takes
the smaller root; by \cref{thm:ch12-ttc} it is negative exactly when
$\pos\cdot\vel < 0$, i.e.\ when the velocity points away from $B$ and the
``collision'' would have happened in the past, so it is reported as no
collision. The wrapper \VoInside forms the relative quantities of
\cref{def:ch12-relative} in line~\ref{alg:ch12-ttc:rel} and compares with the
horizon, which is exactly \cref{def:ch12-truncated}; with $\ttc = \infty$ it
tests membership in the untruncated $\VO_{A|B}$. The test is the same in three
dimensions, where $\pos$ and $\vel$ are 3-vectors (\cref{sec:ch12-3d}).

\begin{algorithm}[htb]
\caption{Sampled velocity selection for one agent}
\label{alg:ch12-choose}
\KwIn{agent $A$ (position, radius, speed limit $v_{\max}$, goal), the other agents $B_1,\dots,B_k$ with their current positions, radii and velocities, horizon $\ttc$, penalty weight $\kappa$}
\KwOut{a new velocity $\vel_A$}
\Fn{\VoChoose{$A$, $\set{B_j}$, $\ttc$}}{
  $\vel^{\mathrm{pref}} \gets$ \VoPreferred{$A$}\tcp*{towards the goal at nominal speed}
  $S \gets$ \VoSamples{$v_{\max}$} $\cup\, \set{\vel^{\mathrm{pref}}}$\label{alg:ch12-choose:samples}\;
  \ForEach{$\vel \in S$}{
    $t_{\min}(\vel) \gets \min_j$ \VoTimeToCollision{$\pos_{B_j}-\pos_A$, $\vel-\vel_{B_j}$, $r_A+r_{B_j}$}\label{alg:ch12-choose:tmin}\;
  }
  $F \gets \set{\vel \in S : t_{\min}(\vel) > \ttc}$\label{alg:ch12-choose:feasible}\;
  \If{$F \ne \emptyset$}{
    \KwRet{$\argmin_{\vel \in F} \norm{\vel - \vel^{\mathrm{pref}}}$}\label{alg:ch12-choose:best}\;
  }
  \KwRet{$\argmin_{\vel \in S} \norm{\vel - \vel^{\mathrm{pref}}} + \kappa / t_{\min}(\vel)$}\label{alg:ch12-choose:fallback}\;
}
\end{algorithm}

\paragraph{Walkthrough of \cref{alg:ch12-choose}.}
Line~\ref{alg:ch12-choose:samples} builds the candidate set: a fixed pattern of
velocities that covers the reachable disc, in the book's code ten concentric
rings of $72$ directions plus the zero velocity ($721$ samples), together with
the preferred velocity itself so that an unobstructed agent flies exactly
where it wants. Line~\ref{alg:ch12-choose:tmin} evaluates every candidate
against every obstacle; the smallest time to collision over the obstacles is
the time to the first collision, and comparing it with $\ttc$ in
line~\ref{alg:ch12-choose:feasible} is membership in the union of
\cref{eq:ch12-feasible}. Line~\ref{alg:ch12-choose:best} is the cost
\cref{eq:ch12-cost} restricted to the samples. If no sample is feasible,
line~\ref{alg:ch12-choose:fallback} chooses the least dangerous one: the
penalty $\kappa/t_{\min}$ is small for a collision far away and huge for one
that is imminent, so the agent buys time. The whole function costs
$\bigO{\abs{S}\,k}$ evaluations of \cref{alg:ch12-ttc}, a few dot products
each; with NumPy the loop over $S$ is one vectorised expression.

\begin{algorithm}[htb]
\caption{Synchronous velocity-obstacle simulation}
\label{alg:ch12-simulate}
\KwIn{agents $A_1,\dots,A_n$ with positions, velocities, radii and goals; time step $\dt$; horizon $\ttc$; number of steps $T$}
\KwOut{the trajectories of all agents}
\Fn{\VoSimulate{$\set{A_i}$, $\dt$, $\ttc$, $T$}}{
  \For{$k \gets 1$ \KwTo $T$}{
    \ForEach{agent $A_i$}{
      \lIf{$A_i$ has reached its goal}{$\vel_i' \gets \vect{0}$}
      \lElseIf{$A_i$ avoids}{$\vel_i' \gets$ \VoChoose{$A_i$, $\set{A_j : j \ne i}$, $\ttc$}\label{alg:ch12-simulate:choose}}
      \lElse{$\vel_i' \gets$ \VoPreferred{$A_i$}\tcp*[f]{a non-cooperative agent}}
    }
    \ForEach{agent $A_i$}{
      $\vel_i \gets \vel_i'$; $\pos_i \gets \pos_i + \dt\,\vel_i$\label{alg:ch12-simulate:move}\;
    }
    \lIf{all agents have reached their goals}{\KwBreak}
  }
  \KwRet{the recorded positions and velocities}\;
}
\end{algorithm}

\paragraph{Walkthrough of \cref{alg:ch12-simulate}.}
The simulation is \emph{synchronous}: in line~\ref{alg:ch12-simulate:choose}
every agent chooses its new velocity from the velocities that the others
\emph{currently} fly, and only then, in line~\ref{alg:ch12-simulate:move}, are
all velocities and positions updated together. This mirrors what happens on
real vehicles, which observe each other's motion and decide at the same time,
and it is the source of the oscillation studied in
\cref{sec:ch12-oscillation}. Agents that do not avoid simply follow their
preferred velocity; they model the non-cooperative intruder. The Euler step in
line~\ref{alg:ch12-simulate:move} is exact for the single-integrator model,
because the velocity is constant during the step. The time step must be
shorter than the horizon, $\dt \le \ttc$, so that a velocity chosen to be
safe for $\ttc$ seconds is safe until the next decision
(\cref{thm:ch12-safety}).

% ---------------------------------------------------------------------
\section{A worked example}
\label{sec:ch12-example}

\begin{example}[Two drones on a crossing course]
\label{ex:ch12-crossing}
Drone $A$ starts at $\pos_A = (-5, 0)$ and wants to reach $(5, 0)$ at its
nominal speed of $1\,\mathrm{m/s}$; it can fly up to $v_{\max} = 1.5\,\mathrm{m/s}$
and applies \cref{alg:ch12-choose} with $\ttc = 5\,\mathrm{s}$ every
$\dt = 0.1\,\mathrm{s}$. Drone $B$ starts at $\pos_B = (0,-5)$, flies north at
$\vel_B = (0, 0.9)$ towards $(0, 5)$ and does not react to $A$. Both have radius
$0.5\,\mathrm{m}$, so $R = 1$. This is \code{crossing\_scenario()} in
\code{code/ch12\_velocity\_obstacles.py}; the function \code{worked\_example()}
prints every number quoted below.
\end{example}

\paragraph{The velocity obstacle at $t=0$.}
The relative position is $\pos_{\mathrm{rel}} = (5,-5)$, at distance
$d = \sqrt{50} = 7.071$ in the direction $-45^\circ$, and the half-angle is
$\theta = \arcsin(1/7.071) = 8.13^\circ$. The two legs of the cone therefore
point at $-36.87^\circ$ and $-53.13^\circ$, which are the directions
$(0.8,-0.6)$ and $(0.6,-0.8)$: the tangent points of the circle of radius $1$
around $(5,-5)$ seen from the origin form $3$--$4$--$5$ triangles. The apex is
$\vel_B = (0, 0.9)$. The preferred velocity of $A$ is
$\vel^{\mathrm{pref}} = (1, 0)$, so the relative velocity is $(1,-0.9)$, which
we already tested in \cref{sec:ch12-ttc}: it collides at
$t_c = 4.559\,\mathrm{s} < \ttc$, so $\vel^{\mathrm{pref}} \in \VO^{\ttc}_{A|B}$.
Indeed, in \cref{fig:ch12-geometry}(b) the point $(1,0)$ lies inside the
truncation disc of centre $(1,-0.1)$ and radius $0.2$, at distance $0.1$ from
its centre.

\paragraph{The choice.}
The nearest point of the boundary of $\VO^{\ttc}_{A|B}$ is straight above
$\vel^{\mathrm{pref}}$ on the truncation arc, $\vel^* = (1, 0.1)$, at distance
$0.1$; its time to collision is exactly $\ttc = 5\,\mathrm{s}$ (the function
\code{closest\_boundary\_point} computes it). The sampled algorithm cannot
return this point because it is not a sample. Of its $722$ candidates, $695$
are feasible, and the feasible one closest to $(1,0)$ is
$\vel_A' = (0.897, 0.078)$, the sample at speed $0.9$ and heading $5^\circ$,
at distance $0.130$ from $\vel^{\mathrm{pref}}$ with $t_c = 5.03\,\mathrm{s}$.
The sampled choice is a little worse than the exact one and, in this case,
slower rather than faster; both are legitimate: a velocity just outside the
truncated cone is safe for the next five seconds, and the decision is
revisited in $0.1\,\mathrm{s}$.

\begin{figure}[tb]
  \centering
  \inputfigure{ch12/example}
  \caption{\Cref{ex:ch12-crossing} simulated with \cref{alg:ch12-simulate}.
  (a)~Trajectories; dots mark whole seconds, the dashed line is $A$'s route
  without avoidance, and the two circles are the discs at $t = 5\,\mathrm{s}$,
  where they come closest. $A$'s detour is barely visible: it passes in front
  of $B$ by flying a little faster and a little to the left. (b)~The velocity
  components of $A$ over time: it speeds up to $1.2\,\mathrm{m/s}$ between
  $t=1$ and $t=3.5\,\mathrm{s}$ and returns to its preferred velocity once
  $B$ has passed. (c)~Centre distance over time with and without avoidance;
  the dashed line is $R = 1$. Without avoidance the discs overlap from
  $t=4.6$ to $5.9\,\mathrm{s}$; with VO the distance never drops below
  $1.001$.}
  \label{fig:ch12-example}
\end{figure}

\paragraph{The run.}
\Cref{tab:ch12-trace} traces the simulation and \cref{fig:ch12-example} plots
it. In the first step $A$ turns $5^\circ$ left and slows to $0.9\,\mathrm{m/s}$.
One step later the relative geometry has changed slightly and the closest
feasible sample is $(1.034, 0.182)$, speed $1.05$ at $10^\circ$. As $B$
approaches, $d$ shrinks, the half-angle $\theta$ grows and the wedge swallows
the candidates near $\vel^{\mathrm{pref}}$ one after another: the number of
feasible samples falls from $695$ at $t=0$ to $122$ at $t = 5\,\mathrm{s}$. From
$t = 1\,\mathrm{s}$ the closest feasible sample is $(1.2, 0)$: fly straight but
faster. This is the typical VO solution to a crossing: it is cheaper to pass in
front of a slower vehicle than to wait for it, and it turns the collision into
a miss because the relative velocity $(1.2,-0.9)$ passes the disc on the far
side (its discriminant is negative). Between $t=4$ and $t=5\,\mathrm{s}$ the
agent adds a small northward component, at $t = 5.0\,\mathrm{s}$ the centres are
$1.001\,\mathrm{m}$ apart, and by $t=5.5\,\mathrm{s}$ the preferred velocity is
feasible again ($t_c = \infty$): $A$ heads for its goal, arriving after a path
of length $10.027\,\mathrm{m}$ instead of $10$. Without avoidance the centres
come within $0.377\,\mathrm{m}$ at $t = 5.2\,\mathrm{s}$; the continuous minimum
is $0.372\,\mathrm{m}$ at the time of closest approach
$t^* = \pos\cdot\vel/(\vel\cdot\vel) = 9.5/1.81 = 5.25\,\mathrm{s}$.

\begin{table}[tb]
  \centering\small
  \caption{Trace of \cref{alg:ch12-choose} for agent $A$ in
  \cref{ex:ch12-crossing}. $t_c(\vel^{\mathrm{pref}})$ is the time to collision
  of the preferred velocity, ``feasible'' the number of feasible candidates
  among the $722$, and $t_c(\vel_A)$ the time to collision of the chosen
  velocity ($\infty$: the ray misses the disc).}
  \label{tab:ch12-trace}
  \begin{tabular}{@{}rrrrrrr@{}}
  \toprule
  $t$ (s) & $\pos_A$ & $\vel^{\mathrm{pref}}-\vel_B$ & $t_c(\vel^{\mathrm{pref}})$ & feasible & chosen $\vel_A$ & $t_c(\vel_A)$\\
  \midrule
  0.0 & $(-5.000, 0.000)$ & $(1.000,-0.900)$ & 4.559 & 695 & $(0.897, 0.078)$ & 5.03\\
  0.1 & $(-4.910, 0.008)$ & $(1.000,-0.901)$ & 4.465 & 694 & $(1.034, 0.182)$ & $\infty$\\
  0.5 & $(-4.497, 0.081)$ & $(1.000,-0.909)$ & 4.091 & 688 & $(1.034, 0.182)$ & $\infty$\\
  1.0 & $(-3.963, 0.154)$ & $(1.000,-0.917)$ & 3.616 & 679 & $(1.200, 0.000)$ & $\infty$\\
  2.0 & $(-2.763, 0.154)$ & $(1.000,-0.920)$ & 2.578 & 642 & $(1.200, 0.000)$ & $\infty$\\
  3.0 & $(-1.563, 0.154)$ & $(1.000,-0.923)$ & 1.574 & 559 & $(1.200, 0.000)$ & $\infty$\\
  4.0 & $(-0.409, 0.181)$ & $(0.999,-0.933)$ & 0.655 & 299 & $(1.046, 0.092)$ & $\infty$\\
  4.5 & $(0.114, 0.227)$ & $(0.999,-0.946)$ & 0.266 & 170 & $(1.046, 0.092)$ & $\infty$\\
  5.0 & $(0.637, 0.273)$ & $(0.998,-0.962)$ & 0.013 & 122 & $(1.050, 0.000)$ & $\infty$\\
  6.0 & $(1.640, 0.215)$ & $(0.998,-0.964)$ & $\infty$ & 657 & $(0.998,-0.064)$ & $\infty$\\
  \bottomrule
  \end{tabular}
\end{table}

Two details of the trace deserve attention. At $t = 5.0\,\mathrm{s}$ the
preferred velocity has $t_c = 0.013\,\mathrm{s}$: the discs are about to touch,
and $\vel^{\mathrm{pref}}$ would push $A$ into $B$ right now. The chosen
velocity $(1.05, 0)$ is nevertheless safe with $t_c = \infty$, because $B$ is
at that moment slightly south of $A$ and any velocity that does not point at
it misses. This is the untruncated-versus-truncated question in miniature: the
wedge is almost $90^\circ$ wide ($d \approx R$), yet more than a hundred of
the samples still lie outside it. Second, the velocity at $t = 6\,\mathrm{s}$
is $(0.998,-0.064)$, which is not on the sample grid: it is the preferred
velocity itself, now pointing slightly south from $(1.64, 0.22)$ towards the
goal, which is why line~\ref{alg:ch12-choose:samples} adds it to the candidates.

% ---------------------------------------------------------------------
\section{Properties: safety, feasibility, complexity}
\label{sec:ch12-properties}

Velocity obstacles are a rule, not a search, so the questions are not
completeness and optimality but: when is the chosen velocity safe, when does a
safe velocity exist, and what does the rule not promise.

\subsection{Safety of one decision}

\begin{theorem}[One-step safety]
\label{thm:ch12-safety}
Let $\norm{\pos_B - \pos_A} > R$ at time $0$, let $A$ choose a velocity
$\vel_A \notin \VO^{\ttc}_{A|B}(\vel_B)$, and let both agents hold their
velocities during $[0,\ttc]$. Then $\norm{\pos_B(t) - \pos_A(t)} > R$ for all
$t \in [0, \ttc]$.
\end{theorem}
\begin{proof}
By \cref{def:ch12-truncated}, $\vel_A \notin \VO^{\ttc}_{A|B}$ means
$t_c(\vel_A - \vel_B) > \ttc$ (possibly $\infty$). By \cref{def:ch12-ttc},
$t_c$ is the first time at which $\norm{\pos_{\mathrm{rel}} - t\vel_{\mathrm{rel}}}
\le R$, and by \cref{eq:ch12-relmotion} this norm is the centre distance at time
$t$ as long as the velocities are held. Hence the distance exceeds $R$ on
$[0, t_c) \supseteq [0, \ttc]$. \qedhere
\end{proof}

\begin{corollary}[A constant-velocity intruder is never hit]
\label{cor:ch12-intruder}
Let $B$ fly with a constant velocity and let $A$ apply \cref{alg:ch12-choose}
at times $0, \dt, 2\dt, \dots$ with $\dt \le \ttc$. If at every decision the
feasible set is non-empty, $A$ and $B$ never collide.
\end{corollary}
\begin{proof}
At each decision time $k\dt$ the agents are apart (by induction), the chosen
velocity is outside $\VO^{\ttc}_{A|B}$, and by \cref{thm:ch12-safety} the
distance stays above $R$ until $k\dt + \ttc \ge (k+1)\dt$, when the next
decision is made. \qedhere
\end{proof}

The corollary is the reason velocity obstacles are the right tool against a
non-cooperative object: the guarantee needs nothing from $B$ except that its
velocity is known and does not change. Every one of the three assumptions has
a practical reading. If $B$'s velocity is \emph{estimated}, from a Kalman
filter for example (\cref{ch:ch18}), the estimation error must be absorbed by
inflating $R$ or by a shorter horizon. If $B$'s velocity \emph{changes}, the
guarantee holds only until it does; the next decision sees the new velocity,
and the horizon must be long enough for $A$ to react. And if $A$ cannot change
its velocity instantaneously, the decision must leave a margin for the time it
takes to reach the new velocity (\cref{sec:ch12-variants}).

\subsection{When no safe velocity exists}

\begin{proposition}[Empty feasible set]
\label[proposition]{thm:ch12-empty}
Let $B$ fly straight at $A$, $\vel_B = -s\,\pos_{\mathrm{rel}}/d$ with speed
$s > v_{\max}$. Then every reachable velocity of $A$ lies in the untruncated
$\VO_{A|B}(\vel_B)$ if and only if $s \ge v_{\max}\,d/R$.
\end{proposition}
\begin{proof}
The apex $\vel_B$ lies on the axis of the wedge at distance $s$ from the
origin, on the side away from $\pos_{\mathrm{rel}}$, so the reachable disc
$\disc{\vect{0}}{v_{\max}}$ is centred on the axis at distance $s > v_{\max}$
from the apex. Seen from the apex it subtends the half-angle
$\arcsin(v_{\max}/s)$, by the same right triangle as in \cref{thm:ch12-cone}.
The disc lies inside the wedge of half-angle $\theta = \arcsin(R/d)$ if and
only if $\arcsin(v_{\max}/s) \le \theta$, i.e.\ $v_{\max}/s \le R/d$.
\qedhere
\end{proof}

In \cref{ex:ch12-crossing} this threshold is $1.5 \cdot 7.07/1 = 10.6\,\mathrm{m/s}$:
an intruder that fast, that close and aimed at $A$ cannot be dodged by any
velocity $A$ can fly. Truncation does not help here, because at that speed
the collision is less than a second away. The fallback in
line~\ref{alg:ch12-choose:fallback} of \cref{alg:ch12-choose} then chooses the
velocity that postpones the collision most, which is the best a purely local
rule can do. In dense traffic the feasible set is also empty, more often and
less dramatically, when many wedges overlap; the experiment of
\cref{sec:ch12-oscillation} counts these events.

\subsection{Complexity and what is not guaranteed}

One evaluation of \cref{alg:ch12-ttc} costs three dot products. With $m$
candidate velocities and $k$ obstacles, \cref{alg:ch12-choose} costs
$\bigO{mk}$, and one step of \cref{alg:ch12-simulate} with $n$ avoiding agents
costs $\bigO{n^2 m}$; for $m = 722$ this is about $5 \cdot 10^5$ evaluations
per step for $n = 8$, which NumPy handles in a few milliseconds. Restricting
each agent to its $k$ nearest neighbours, as every real implementation does,
makes the step linear in $n$.

What the rule does \emph{not} guarantee is as important as what it does.
\begin{itemize}
  \item \textbf{No completeness or optimality.} The rule is greedy over one
  step. A velocity that is safe now can lead to a position from which every
  velocity is unsafe later, for instance when an agent flies into the
  narrowing gap between two converging obstacles. Fiorini and Shiller
  addressed this with a tree search over sequences of avoidance
  manoeuvres~\cite{fiorini1998motion}; the one-step rule is what practical
  systems use, backed by a global planner (\cref{ch:ch24}).
  \item \textbf{No progress guarantee.} Nothing forces the agent towards its
  goal. Two agents that must pass each other in a corridor narrower than
  $2R$ stop and stay stopped; the deadlock is a task for the replanning
  layer, not for the local rule.
  \item \textbf{No guarantee among cooperating agents.} \Cref{thm:ch12-safety}
  assumes that $B$ holds its velocity. When $B$ is itself avoiding $A$, the
  assumption fails at every step, and safety and smoothness both suffer:
  this is the subject of the next section.
  \item \textbf{Sampling error.} The sampled minimiser can be off the exact
  one by the spacing of the sample grid, and a thin feasible corridor between
  two wedges can be missed entirely, so that the algorithm declares the set
  empty although it is not (\cref{sec:ch12-implementation}).
\end{itemize}

% ---------------------------------------------------------------------
\section{The oscillation problem}
\label{sec:ch12-oscillation}

\subsection{Two agents avoiding each other at the same time}

Let two identical agents fly head-on towards each other's start positions,
$A$ from $(-5,0)$ and $B$ from $(5, 0.2)$, both at speed $1$ with $R = 1$, and
let \emph{both} apply \cref{alg:ch12-choose} at every step. At $t = 0$ each
sees the other coming straight at it and steps aside: $A$ chooses the sample
$(1.034, -0.182)$, $B$ its mirror image $(-1.034, 0.182)$. Now consider the
next decision. $A$ assumes that $B$ will keep its new, sideways velocity. With
$B$ moving up and to the left, the relative velocity of the \emph{preferred}
velocity $(1,0)$ now clears the disc, so $A$ returns to $(1,0)$, and $B$, by
symmetry, returns to $(-1,0)$. One step later both are heading straight at
each other again, step aside again, and so on. \Cref{fig:ch12-oscillation}
shows the result: the lateral velocity of $A$ changes sign $51$ times in
$100$ steps, and although the minimum distance stays at $1.002 \ge R$ the
velocity command chatters between $\pm0.18$ and $0$ at every step. A real
drone cannot follow such a command, and the manoeuvre it executes is not the
one that was analysed.

\begin{figure}[tb]
  \centering
  \inputfigure{ch12/oscillation}
  \caption{Oscillation. Two agents swap places head-on with a lateral offset
  of $0.2$. (a)~Trajectories when both apply the velocity-obstacle rule
  (solid) and when only $A$ applies it (dashed): with two avoiders the drift
  is shared and the paths zigzag. (b)~The lateral velocity of $A$: with two
  avoiders it flips sign at almost every step, because each agent reacts to
  the other's sidestep by cancelling its own; with one avoider it changes
  sign once. The vertical axis of (a) is stretched.}
  \label{fig:ch12-oscillation}
\end{figure}

The cause is the assumption behind \cref{def:ch12-vo}: $B$ keeps its velocity.
When $B$ is an avoider too, its velocity at the next step is a \emph{reaction}
to $A$'s choice, and the two reactions cancel. The remedy is to let each agent
take only half of the responsibility for the avoidance and to assume that the
other takes the other half: this is the reciprocal velocity obstacle of van
den Berg, Lin and Manocha~\cite{vandenberg2008rvo}, whose apex is at the
average $(\vel_A + \vel_B)/2$ instead of $\vel_B$, and its refinements HRVO
\cite{snape2011hrvo} and \orca~\cite{vandenberg2011orca}, which turns the
choice into a linear program. \Cref{ch:ch13} derives them from the geometry of
this chapter. Note that reciprocity is a property of the \emph{pair}: against
an intruder that does not react, the plain velocity obstacle with full
responsibility is the correct construction, and using a reciprocal one would
avoid only half of the collision.

\subsection{An experiment on crossing agents}
\label{sec:ch12-experiment}

The script \code{code/figures/gen\_ch12\_sim.py} runs two families of
scenarios with the simulator of \cref{alg:ch12-simulate} ($\dt = 0.1$,
$\ttc = 5$, $R = 1$, nominal speed $1$, $v_{\max} = 1.5$). In the
\emph{stream} scenario one avoiding agent flies $12\,\mathrm{m}$ east across a
picket line of $k$ non-cooperative intruders that fly north at speeds
$0.8$, $1.1$ and $1.4\,\mathrm{m/s}$, each timed to hit the agent exactly where
its nominal path crosses theirs. In the \emph{circle} scenario $n$ agents start
evenly spaced on a circle of radius $6$ and fly to the antipodal points, so
every pair is on a crossing course through the centre, and \emph{all} of them
apply the velocity-obstacle rule (the start positions are perturbed by a
fixed-seed noise of standard deviation $0.05$ to break the symmetry).
\Cref{fig:ch12-crossing} and \cref{tab:ch12-crossing} report the minimum
centre distance over the run and the path length relative to the straight
line.

\begin{figure}[tb]
  \centering
  \inputfigure{ch12/crossing}
  \caption{No avoidance versus velocity obstacles for $n = 2, \dots, 8$
  crossing agents. (a)~Minimum centre distance; the dashed line is $R = 1$.
  Without avoidance every scenario collides. One VO agent among
  non-cooperative intruders (stream) keeps the distance at exactly $R$ for
  every $k$, as \cref{cor:ch12-intruder} promises. When all agents apply VO
  simultaneously (circle) the distance dips below $R$ from $n=5$ on.
  (b)~Path length relative to the straight line: the price of avoidance is
  $1$--$2\,\%$ in the stream and up to $14\,\%$ in the crowded circle.}
  \label{fig:ch12-crossing}
\end{figure}

\begin{table}[tb]
  \centering\small
  \caption{The experiment of \cref{fig:ch12-crossing}. ``Pairs $<R$'' counts
  the pairs of agents whose centre distance dropped below $R$ at some step,
  ``empty'' the number of decisions at which no sampled velocity was feasible.}
  \label{tab:ch12-crossing}
  \begin{tabular}{@{}r rrr rrrrr@{}}
  \toprule
  & \multicolumn{3}{c}{stream: one VO agent, $k=n$ intruders} & \multicolumn{5}{c}{circle: $n$ VO agents}\\
  \cmidrule(lr){2-4}\cmidrule(lr){5-9}
  $n$ & min dist none & min dist VO & path ratio & min dist none & min dist VO & path ratio & pairs $<R$ & empty\\
  \midrule
  2 & 0.000 & 1.000 & 1.014 & 0.038 & 1.009 & 1.008 & 0 & 0\\
  3 & 0.000 & 1.000 & 1.019 & 0.016 & 0.990 & 1.010 & 1 & 0\\
  4 & 0.000 & 1.000 & 1.019 & 0.021 & 1.001 & 1.014 & 0 & 0\\
  5 & 0.000 & 1.000 & 1.019 & 0.024 & 0.976 & 1.048 & 4 & 1\\
  6 & 0.000 & 1.000 & 1.020 & 0.007 & 0.914 & 1.142 & 6 & 5\\
  7 & 0.000 & 1.000 & 1.018 & 0.003 & 0.936 & 1.061 & 4 & 4\\
  8 & 0.000 & 1.000 & 1.023 & 0.005 & 0.954 & 1.115 & 5 & 1\\
  \bottomrule
  \end{tabular}
\end{table}

The two halves of the table tell the two stories of this chapter. Against
non-cooperative intruders, the velocity-obstacle agent never comes closer
than $R$ (the recorded minimum is $1.0003$), whatever the number of intruders,
and pays for it with $1.4$--$2.3\,\%$ extra path; it even arrives earlier than
the $12\,\mathrm{s}$ of the straight flight, in $11.3$--$11.5\,\mathrm{s}$,
because its avoidance manoeuvre is mostly ``speed up and pass in front''.
Among agents that all apply the rule simultaneously, the picture degrades
with density: up to $n = 4$ the runs are clean, but from $n = 5$ on the
minimum distance falls below $R$ (to $0.914$ for $n=6$, an overlap of
$9\,\%$ of the combined radius), several pairs graze each other, decisions
with an empty feasible set appear, and the paths lengthen by up to $14\,\%$
as agents oscillate around the crowded centre. None of these are head-on
crashes, and $n = 8$ agents on a circle of radius $6$ with $R = 1$ is a
crowded scene; but the guarantee of \cref{cor:ch12-intruder} is gone, because
its assumption is violated at every step. \Cref{ch:ch13} repeats this
experiment with RVO and \orca.

% ---------------------------------------------------------------------
\section{Variants and extensions}
\label{sec:ch12-variants}

\paragraph{Static and polygonal obstacles.}
A static obstacle has $\vel_B = \vect{0}$, so its velocity obstacle is the
collision cone itself with the apex at the origin, and its truncated version
is the cone beyond the disc $\disc{\pos_{\mathrm{rel}}/\ttc}{R/\ttc}$. An obstacle
that is not a disc, such as a building or a wall, is first inflated by $r_A$
(\cref{ch:ch02}) and then treated point by point: a velocity is dangerous if
its ray hits any point of the inflated obstacle, so the cone is the union of
the cones of all its points. For a convex obstacle this union is the wedge
between the two extreme tangent rays from $A$ (\cref{exr:ch12-static}). The
truncated version is again the union of the scaled copies of the obstacle,
$\bigcup_{t \le \ttc} \mathcal{O}/t$, which for a polygon is easy to draw and
easy to sample. Truncation matters most for static obstacles: without it a
long wall forbids every velocity with any component towards it, however far
away the wall is.

\paragraph{Obstacles that do not fly straight.}
The velocity obstacle assumes that $B$ keeps its velocity. If the prediction
layer (\cref{ch:ch20}) supplies a predicted trajectory $\hat{\pos}_B(t)$
instead, \cref{thm:ch12-truncated} generalises without change of argument: $A$
flying straight with velocity $\vel$ is within $R$ of $B$ at time $t$ if and
only if $\norm{\hat{\pos}_B(t) - \pos_A - t\vel} \le R$, i.e.\
$\vel \in \disc{(\hat{\pos}_B(t) - \pos_A)/t}{R/t}$, so the forbidden set is
the union of these discs over $t \le \ttc$. It is no longer a cone, but the
sampled \cref{alg:ch12-choose} handles it with the membership test replaced
by a scan over the predicted times. This is the non-linear velocity obstacle.
Uncertainty in the prediction is absorbed by inflating $R$ with the predicted
standard deviation, as the prediction-aware constraints of \cref{ch:ch24} do.

\paragraph{Reciprocal variants.}
For agents that all avoid each other, the reciprocal velocity obstacle
\cite{vandenberg2008rvo} shifts the apex to $(\vel_A + \vel_B)/2$, the hybrid
reciprocal velocity obstacle \cite{snape2011hrvo} combines one leg of RVO with
one leg of VO to fix the side on which the agents pass, and \orca
\cite{vandenberg2011orca} replaces each wedge by a half-plane so that the
selection becomes a small linear program. All three are derived in
\cref{ch:ch13} from \cref{def:ch12-vo} and \cref{thm:ch12-truncated}.

\paragraph{Dynamics.}
\Cref{def:ch12-feasible} takes every speed below $v_{\max}$ to be reachable at
once. With an acceleration limit $a_{\max}$ the velocities reachable within one
step form the disc $\disc{\vel_A}{a_{\max}\dt}$ intersected with
$\disc{\vect{0}}{v_{\max}}$, the dynamic window of \cref{ch:ch14}; the sampled
algorithm accepts any candidate set, so it is enough to sample that window
instead. What does not carry over is the exactness of \cref{thm:ch12-safety}:
while the agent accelerates towards its new velocity it flies neither the old
nor the new one. A margin on $R$ or on the horizon, of the order of the time
$\norm{\vel_A' - \vel_A}/a_{\max}$ needed to reach the new velocity, restores
safety in practice.

\paragraph{Exact selection.}
With a single obstacle, the minimiser of \cref{eq:ch12-cost} is the projection
of $\vel^{\mathrm{pref}}$ onto the boundary of $\VO^{\ttc}_{A|B}$: onto one of
the two legs beyond their tangency points or onto the truncation arc,
whichever is closest. The function \code{closest\_boundary\_point} in the
chapter's code does this in a few lines and returned $\vel^* = (1, 0.1)$ in
\cref{ex:ch12-crossing}. With several obstacles the projection onto one wedge
may land inside another, and the exact answer is the closest point of a
region bounded by several arcs and legs; \orca avoids this by using one
half-plane per obstacle.

% ---------------------------------------------------------------------
\section{The cone in three dimensions}
\label{sec:ch12-3d}

Drones fly in three dimensions, and the third dimension is often the cheapest
way out of a conflict: climbing over an intruder can be a smaller manoeuvre
than swerving around it. Nothing in the derivation above used the plane. Let
$\pos_A, \vel_A, \pos_B, \vel_B \in \R^3$ and let the agents be spheres of
radii $r_A$, $r_B$. The relative quantities of \cref{def:ch12-relative}, the
collision condition \cref{eq:ch12-collision} and the quadratic
\cref{eq:ch12-quadratic} are unchanged, with 3-vectors in the dot products,
because the Minkowski sum of two balls is again a ball. The ray in direction
$\vel$ meets the ball $\disc{\pos_{\mathrm{rel}}}{R}$ if and only if the
distance from the centre to the ray, $d\sin\varphi$, is at most $R$, exactly as
in the proof of \cref{thm:ch12-cone}, and the angle $\varphi$ between $\vel$
and $\pos_{\mathrm{rel}}$ is now the angle between two vectors in space.

\begin{proposition}[The 3D collision cone]
\label[proposition]{thm:ch12-cone3d}
For spheres at distance $d > R$, the collision cone $CC_{A|B} \subset \R^3$ is
the right circular cone\index{velocity obstacle!in three dimensions} with
apex at the origin, axis $\pos_{\mathrm{rel}}$ and half-angle
$\theta = \arcsin(R/d)$; the velocity obstacle is its translate by $\vel_B$,
and the truncated velocity obstacle is the cone with its tip cut off by the
ball $\disc{\vel_B + \pos_{\mathrm{rel}}/\ttc}{R/\ttc}$, which is inscribed in
the cone. Every plane through the axis cuts the cone in the planar wedge of
\cref{thm:ch12-cone}.
\end{proposition}

\begin{figure}[tb]
  \centering
  \inputfigure{ch12/cone3d}
  \caption{The collision cone in three dimensions: a right circular cone
  with apex at $A$, axis towards the centre of the sphere of radius
  $R = r_A + r_B$ around $B$ and half-angle $\theta = \arcsin(R/d)$. The cone
  touches the sphere along a circle. Any plane through the axis shows the
  planar picture of \cref{fig:ch12-idea}; a relative velocity $\vel$ is
  dangerous exactly when it lies inside the cone.}
  \label{fig:ch12-cone3d}
\end{figure}

\Cref{fig:ch12-cone3d} shows the cone. The proof is the one of
\cref{thm:ch12-cone} and \cref{thm:ch12-truncated} read with 3-vectors; the
only new fact is that the set of directions at angle $\le \theta$ from a fixed
axis in $\R^3$ is a circular cone rather than a wedge. The tangent points form
a circle at distance $(d^2-R^2)/d$ from the apex along the axis, of radius
$R\sqrt{d^2-R^2}/d$. For the implementation this means that \cref{alg:ch12-ttc}
works verbatim in 3D, which the chapter's code exploits: \code{time\_to\_collision}
accepts 3-vectors, and its self-test checks the cone for
$\pos_{\mathrm{rel}} = (3,4,12)$, $d = 13$ and $R = 1$ by probing directions
just inside and just outside the half-angle $\arcsin(1/13)$. What changes is
the candidate set of \cref{alg:ch12-choose}: the reachable ball must be
sampled in speed and in two angles, for instance with a few hundred
well-spread directions on the unit sphere times a handful of speeds, and a
drone will usually weight vertical deviations differently from horizontal ones
in the cost \cref{eq:ch12-cost}, because climbing costs energy and altitude
bands may be reserved. The reciprocal constructions of \cref{ch:ch13} have the
same 3D form.

% ---------------------------------------------------------------------
\section{Implementation notes}
\label{sec:ch12-implementation}

The file \code{code/ch12\_velocity\_obstacles.py} implements everything in this
chapter in about 350 lines: the plane-geometry helpers copied from
\code{code/ch02\_toolbox.py}, the time to collision and membership test
(\cref{lst:ch12-ttc}), a \code{VelocityObstacle} class that exposes the apex,
the tangent directions, the truncation disc and the exact projection onto the
boundary, the sampled selection (\cref{lst:ch12-choose}), the simulator of
\cref{alg:ch12-simulate} and the scenarios of the figures. Its self-test runs
in well under a second.

\paragraph{Numerical care when agents nearly touch.}
Three situations need explicit handling. When $d \le R$ the discs already
overlap, the constant term of \cref{eq:ch12-quadratic} is non-positive and
every velocity ``collides'' at $t = 0$; treating all of them as forbidden would
freeze the agent inside the other. The code instead forbids the velocities that
approach ($\pos_{\mathrm{rel}}\cdot\vel_{\mathrm{rel}} > 0$) and allows those
that separate, so that an overlap, whether from a bad initial placement or a
violated assumption, resolves itself. When $d$ is only slightly larger than
$R$, the half-angle is close to $90^\circ$ and the discriminant of a grazing
velocity hovers around zero; rounding then decides between ``graze'' and
``miss''. Do not rely on the exact boundary: the exact projection
$\vel^* = (1,0.1)$ of \cref{ex:ch12-crossing} has $t_c = \ttc$ to the last
digit, and a real agent should either inflate $R$ by a margin or demand
$t_c > \ttc + \epsilon$. Finally, $\vel_{\mathrm{rel}} = \vect{0}$ (two
hovering drones, or two drones in formation) makes the quadratic degenerate;
\cref{alg:ch12-ttc} returns $\infty$ for it, which is correct because the
distance does not change.

\paragraph{The horizon.}
$\ttc$ must be longer than the decision interval $\dt$ (\cref{thm:ch12-safety})
and long enough for the manoeuvre that a velocity just outside
$\VO^{\ttc}$ may require at the next decision: a few times the time needed to
fly a distance $R$, plus the time to change velocity under the acceleration
limit. Too short a horizon produces late, sharp manoeuvres and empty feasible
sets; too long a horizon produces the frozen behaviour of the pitfall below.
The chapter's examples use $\ttc = 5\,\mathrm{s}$ for $R = 1\,\mathrm{m}$ and
speeds around $1\,\mathrm{m/s}$; scale it with $R/v$.

\paragraph{Sampling versus solving exactly.}
The default candidate set has $721$ velocities on ten rings of $72$
directions, i.e.\ $5^\circ$ in heading and $0.15\,\mathrm{m/s}$ in speed for
$v_{\max} = 1.5$. Always add the preferred velocity and, to reduce chattering,
the current velocity to the candidates, and consider a finer ring around
$\vel^{\mathrm{pref}}$. With a single obstacle the exact projection of
\code{closest\_boundary\_point} is both cheaper and better; with many
obstacles, sampling is robust and trivially accepts other cost functions and
reachable sets (the dynamic window of \cref{ch:ch14}), while the exact
solution is the linear program of \orca in \cref{ch:ch13}. The test of all
candidates against one obstacle is one vectorised NumPy expression
(\code{ttc\_batch}), which is why the whole self-test, including the
experiments of \cref{sec:ch12-experiment}, runs in about $0.2\,\mathrm{s}$.

\begin{lstlisting}[caption={Time to collision and membership in the truncated velocity obstacle (from \code{code/ch12\_velocity\_obstacles.py}); this is \cref{alg:ch12-ttc}.},label={lst:ch12-ttc}]
def time_to_collision(p_rel, v_rel, radius):
    """First time t >= 0 at which |p_rel - t v_rel| <= radius, or math.inf
    if the discs (spheres) never touch.  Works in 2D and in 3D."""
    p = np.asarray(p_rel, dtype=float)
    v = np.asarray(v_rel, dtype=float)
    pp, pv, vv = p @ p, p @ v, v @ v
    k = pp - radius * radius
    if k <= 0.0:
        return 0.0                        # already overlapping
    if vv == 0.0:
        return math.inf                   # no relative motion
    disc = pv * pv - vv * k               # discriminant / 4
    if disc < 0.0:
        return math.inf                   # the ray misses the disc
    t = (pv - math.sqrt(disc)) / vv       # smaller root = first contact
    return t if t >= 0.0 else math.inf    # the disc lies behind A


def in_velocity_obstacle(p_rel, v_rel, radius, tau=math.inf):
    """True iff the relative velocity v_rel leads to a collision within the
    horizon tau (tau = inf gives the untruncated velocity obstacle)."""
    tc = time_to_collision(p_rel, v_rel, radius)
    return tc < math.inf and tc <= tau
\end{lstlisting}

\Cref{lst:ch12-ttc} is \cref{alg:ch12-ttc} line by line. Note the order of the
guards: the overlap test comes before the test for zero relative velocity,
because two overlapping agents at rest must still be told that they are
colliding. \Cref{lst:ch12-choose} is \cref{alg:ch12-choose}. The candidates
are a NumPy array with one row per velocity; for every other agent the
relative velocities of all candidates are formed at once and passed to the
vectorised \code{ttc\_batch}, and the running minimum over obstacles is the
$t_{\min}$ of line~\ref{alg:ch12-choose:tmin}. The dictionary \code{info}
returns the numbers reported in \cref{tab:ch12-trace}.

\begin{lstlisting}[caption={Sampled velocity selection (from \code{code/ch12\_velocity\_obstacles.py}); this is \cref{alg:ch12-choose}.},label={lst:ch12-choose}]
def choose_velocity(agent, others, tau, dt, samples=None, penalty=1.0):
    """Sampled velocity-obstacle selection for ``agent`` (Algorithm 12.2).

    Every candidate is tested against the truncated VO of every other
    agent (others keep their current velocities).  Among the candidates
    that are outside all VO^tau, the one closest to the preferred velocity
    wins; if there is none, the penalised cost |v - v_pref| + penalty/t_c
    picks the least dangerous candidate.  Returns (velocity, info)."""
    v_pref = agent.preferred_velocity(dt)
    if samples is None:
        samples = velocity_samples(agent.v_max)
    cand = np.vstack([samples, np.asarray([v_pref])])
    tc = np.full(len(cand), np.inf)
    for other in others:
        if other is agent:
            continue
        p_rel = sub(other.position, agent.position)
        v_rel = cand - np.asarray(other.velocity, dtype=float)
        radius = agent.radius + other.radius
        if norm(p_rel) <= radius:            # already overlapping: forbid
            tc_j = np.where(v_rel @ np.asarray(p_rel) > 0.0, 0.0, np.inf)
        else:                                # every closing velocity
            tc_j = ttc_batch(p_rel, v_rel, radius)
        tc = np.minimum(tc, tc_j)
    dist = np.linalg.norm(cand - np.asarray(v_pref), axis=1)
    feasible = (tc > tau) | np.isinf(tc)
    if feasible.any():
        idx = int(np.argmin(np.where(feasible, dist, np.inf)))
    else:
        idx = int(np.argmin(dist + penalty / np.maximum(tc, 1e-9)))
    info = {"v_pref": v_pref, "tc_pref": float(tc[-1]),
            "n_feasible": int(feasible.sum()), "tc": float(tc[idx]),
            "feasible": bool(feasible[idx]), "dist": float(dist[idx])}
    return (float(cand[idx, 0]), float(cand[idx, 1])), info
\end{lstlisting}

\begin{pitfall}[Forgetting to add the radii]
The disc in \cref{def:ch12-collision-cone} has radius $r_A + r_B$, not $r_B$.
Using $B$'s radius alone (or, worse, treating both agents as points) makes the
cone too narrow by the missing $\arcsin$ term, and the agent then grazes or
overlaps the obstacle by up to $r_A$ while its own test reports safety. In a
simulation this shows up as a minimum separation between $r_B$ and $R$; in
the experiment of \cref{tab:ch12-crossing} it would turn every ``$1.000$''
into ``$0.5$''. Add the radii, then add a margin for estimation error and
dynamics.\index{pitfall!forgetting the radii}
\end{pitfall}

\begin{pitfall}[Absolute instead of relative velocity]
The cone is tested with $\vel_A - \vel_B$, and the apex of the velocity
obstacle is at $\vel_B$. Testing $\vel_A$ itself against the cone at the
origin is correct only for static obstacles and silently wrong for moving
ones: the agent then avoids where the obstacle \emph{is}, not where it will
be, and collides with anything that moves across its path. A test that works
against walls but fails against other drones is almost always this
mistake.\index{pitfall!absolute velocity}
\end{pitfall}

\begin{pitfall}[No truncation]
With $\ttc = \infty$ the velocity obstacle forbids every velocity that would
collide at \emph{any} future time. A wall then forbids every velocity with a
component towards it, and a slow drone a hundred metres away blocks a whole
sector of headings; the agent refuses far-future collisions that a later
decision would have handled, creeps, or freezes short of its goal. Truncate
with a horizon (\cref{def:ch12-truncated}) and choose it as described above.
\index{pitfall!no truncation}
\end{pitfall}

\begin{pitfall}[Discretisation of the velocity samples]
A coarse sample set has two failure modes. A thin feasible corridor between
two wedges, which is common when two obstacles pass on either side, may
contain no sample at all, and the algorithm declares the feasible set empty
and falls back to a dangerous velocity although a safe one exists. And when
the best sample changes from step to step by one grid cell, the command
chatters. Sample finely near the preferred and the current velocity, always
include both as candidates, and prefer the exact projection when there is a
single obstacle.\index{pitfall!velocity sampling}
\end{pitfall}

% ---------------------------------------------------------------------
\section{Where this fits in the drone system}
\label{sec:ch12-drone}

\begin{dronebox}
Velocity obstacles and their reciprocal descendants form the \emph{local
safety layer} of the hybrid architecture in \cref{ch:ch24}. The executive
monitors, for every drone, the time to collision \cref{eq:ch12-ttc} against
every tracked object, using the positions and velocities estimated by the
prediction layer (\cref{ch:ch18,ch:ch20}). While the time to collision along
the nominal route is larger than the safety horizon, the drone follows the
route computed by the global planner (\cref{ch:ch09,ch:ch10}). When it drops
below the horizon, the local layer takes over the velocity command: the
preferred velocity is the velocity along the nominal route, the obstacles are
the tracked objects with their estimated velocities, and the radius $R$ is
inflated by a margin that grows with the uncertainty of the estimate. Once no
collision is predicted within the horizon, control returns to the route, and
if the manoeuvre has pushed the drone too far off its route the replanning
layer of \cref{ch:ch05} is invoked.

Which construction the layer uses depends on who the obstacle is. A
non-cooperative intruder, an unknown drone or a bird, is exactly the case of
this chapter: it does not react, so the drone takes the full responsibility
and uses $\VO^{\ttc}$ with the apex at the intruder's estimated velocity
(\cref{cor:ch12-intruder}). Another member of the swarm reacts with the same
rule, and the reciprocal constructions of \cref{ch:ch13} apply. Mixing them
up is a real error: a reciprocal manoeuvre against an intruder avoids only half
of the collision, and a full-responsibility manoeuvre among swarm members
produces the oscillation of \cref{sec:ch12-oscillation}. This chapter is the
first half of Week~6 of the study plan (\cref{ch:appA}); the coding exercise
\cref{exr:ch12-coding} builds the crossing simulation that the second half,
\cref{ch:ch13}, extends with RVO and \orca.
\end{dronebox}

% ---------------------------------------------------------------------
\section{Summary}
\begin{summary}
\begin{itemize}
  \item Local collision avoidance chooses a \emph{velocity} for the next
  instant instead of a path. In the frame of agent $A$, agent $B$ is a disc of
  radius $R = r_A + r_B$ at $\pos_{\mathrm{rel}} = \pos_B - \pos_A$, and the
  relative velocity $\vel_{\mathrm{rel}} = \vel_A - \vel_B$ collides exactly
  when its ray hits that disc.
  \item The collision cone $CC_{A|B}$ is the wedge of half-angle
  $\theta = \arcsin(R/d)$ around $\pos_{\mathrm{rel}}$, bounded by the tangent
  rays to the disc; the velocity obstacle $\VO_{A|B} = \vel_B \oplus CC_{A|B}$
  is the same wedge with its apex at $\vel_B$.
  \item The time to collision is the smaller root of
  $(\vel\cdot\vel)t^2 - 2(\pos\cdot\vel)t + (\pos\cdot\pos - R^2) = 0$; a
  collision exists if and only if the discriminant is non-negative and
  $\pos\cdot\vel > 0$.
  \item The truncated velocity obstacle $\VO^{\ttc}_{A|B}$ keeps only the
  velocities that collide within the horizon $\ttc$; it is the wedge with its
  tip cut off by the inscribed disc $\disc{\vel_B + \pos_{\mathrm{rel}}/\ttc}{R/\ttc}$,
  and it is what makes the method usable near walls and in traffic.
  \item The agent chooses, among the reachable velocities outside the union of
  all truncated velocity obstacles, the one closest to its preferred velocity;
  sampling the velocity disc is the simple and robust way to do so. Static
  obstacles are the case $\vel_B = \vect{0}$, and the construction is the same
  in three dimensions with a circular cone.
  \item A velocity outside $\VO^{\ttc}$ is safe for $\ttc$ seconds if the
  obstacle holds its velocity; hence a constant-velocity intruder is never hit.
  When both agents apply the rule at once, the assumption fails, the commands
  oscillate and, in dense scenes, the separation is no longer guaranteed. That
  is what the reciprocal methods of \cref{ch:ch13} repair.
\end{itemize}
\end{summary}

\paragraph{Further reading.}
The velocity obstacle was introduced by Fiorini and
Shiller~\cite{fiorini1998motion}, whose paper also contains the distinction
between imminent and distant collisions and a tree search over sequences of
avoidance manoeuvres; the geometric background (Minkowski sums,
configuration space) is covered in the textbook of Choset et
al.~\cite{choset2005principles}. The reciprocal velocity obstacle
\cite{vandenberg2008rvo}, the hybrid reciprocal velocity obstacle
\cite{snape2011hrvo} and \orca~\cite{vandenberg2011orca} are the subject of
\cref{ch:ch13}; the dynamic window approach of \cref{ch:ch14} shares the
velocity-space view but replaces the collision cone by a stopping-distance
test.

% ---------------------------------------------------------------------
\section{Exercises}
\label{sec:ch12-exercises}

\begin{exercise}[\difficulty{1}]
\label{exr:ch12-halfangle}
Two drones of radii $r_A = 0.3\,\mathrm{m}$ and $r_B = 0.5\,\mathrm{m}$ are
$4\,\mathrm{m}$ apart. Compute the half-angle of the collision cone. Recompute
it with $r_A$ forgotten (radius $r_B$ only) and with both radii forgotten, and
give the width of the wedge in degrees in each case. At what distance does the
correct half-angle reach $30^\circ$, and what happens to the cone as the
distance approaches $r_A + r_B$?
\end{exercise}

\begin{exercise}[\difficulty{2} Proof]
\label{exr:ch12-tangent}
(a)~Using only the fact that a tangent to a circle is perpendicular to the
radius at the point of contact, prove that the two tangent rays from a point
at distance $d$ to a circle of radius $R < d$ make the angle
$\theta = \arcsin(R/d)$ with the direction to the centre, and that they touch
the circle at distance $\sqrt{d^2 - R^2}$ from the point. (b)~Show that
$CC_{A|B}$ is closed under multiplication by positive scalars and that
$\VO_{A|B}(\vel_B)$ is closed under it if and only if $\vel_B = \vect{0}$.
(c)~Conclude that whether an \emph{untruncated} velocity obstacle contains
$\vel_A$ depends only on the direction of $\vel_A - \vel_B$, not on its
magnitude, and explain why this is no longer true after truncation.
\end{exercise}

\begin{exercise}[\difficulty{2} Proof]
\label{exr:ch12-truncation}
Prove \cref{thm:ch12-truncated} in detail: show that a relative velocity
$\vel$ collides at time $t$ if and only if $\vel \in \disc{\pos_{\mathrm{rel}}/t}{R/t}$,
that each of these discs is tangent to both legs of the cone, and that the
point of the truncated obstacle closest to the apex lies on the axis at
distance $(d - R)/\ttc$, while the tangency points of the truncation disc lie
at distance $\sqrt{d^2 - R^2}/\ttc$ along the legs. Verify both distances in
\cref{fig:ch12-geometry}(b), where $d = \sqrt{50}$, $R = 1$ and $\ttc = 5$.
\end{exercise}

\begin{exercise}[\difficulty{1}]
\label{exr:ch12-ttc}
Let $\pos_{\mathrm{rel}} = (6, 2)$, $R = 2.5$ and $\vel_{\mathrm{rel}} = (3, 0)$.
Write down the quadratic \cref{eq:ch12-quadratic}, compute the time to
collision and the time at which the discs separate again, and compute the
time and distance of closest approach. Repeat for $\vel_{\mathrm{rel}} = (-3, 0)$
and explain the result with \cref{thm:ch12-ttc}. Check both with
\code{time\_to\_collision}.
\end{exercise}

\begin{exercise}[\difficulty{2}]
\label{exr:ch12-empty}
An intruder $B$ flies straight at $A$ with speed $s$ from distance $d$. Derive
the condition of \cref{thm:ch12-empty}, $s \ge v_{\max} d/R$, from the tangent
half-angle of the reachable disc seen from the apex, and evaluate it for
$v_{\max} = 1.5\,\mathrm{m/s}$, $d = 7.07\,\mathrm{m}$, $R = 1\,\mathrm{m}$. For
$s$ just above the threshold, which velocity does the fallback of
\cref{alg:ch12-choose} choose, and why does truncation with $\ttc = 5\,\mathrm{s}$
not create a feasible velocity? Then describe the situation with $s$ well
below the threshold but $d$ small: which velocities are feasible, and what
does this say about the horizon?
\end{exercise}

\begin{exercise}[\difficulty{2}]
\label{exr:ch12-static}
(a)~Show that for a static obstacle the velocity obstacle equals the
collision cone and that the truncation disc has centre $\pos_{\mathrm{rel}}/\ttc$.
(b)~Let $\mathcal{O}$ be a convex obstacle (already inflated by $r_A$) that
does not contain the origin. Show that the set of velocities whose ray meets
$\mathcal{O}$ is the wedge between the two extreme tangent rays from the
origin to $\mathcal{O}$, i.e.\ the union of the collision cones of all points
of $\mathcal{O}$ is itself a cone. (c)~A drone flies parallel to a long
straight wall at distance $2\,\mathrm{m}$ from its centre, with $r_A = 0.5$.
Which velocities does the untruncated velocity obstacle of the wall forbid?
Which ones does the truncated obstacle with $\ttc = 4\,\mathrm{s}$ forbid?
\end{exercise}

\begin{exercise}[\difficulty{3} Coding]
\label{exr:ch12-coding}
This is the first half of the coding exercise of Week~6 of the study plan
(\cref{ch:appA}). Build a 2D simulation of $n$ disc agents on crossing courses
(the circle and stream scenarios of \cref{sec:ch12-experiment}, or your own),
either from scratch following \cref{alg:ch12-ttc,alg:ch12-choose,alg:ch12-simulate}
or by extending \code{code/ch12\_velocity\_obstacles.py}. Compare no avoidance
with velocity obstacles and report, per run, the minimum centre distance, the
path length relative to the straight line, the travel time, the number of
decisions with an empty feasible set and the number of sign changes of the
lateral velocity (a measure of oscillation). Then (a)~vary the horizon
$\ttc \in \set{1, 2, 5, 10, \infty}$ and find the values for which the agents
freeze or collide; (b)~vary the sample resolution ($36$, $72$ and $180$
directions) and observe the chattering; (c)~add a static obstacle as a wall of
discs and check the frozen behaviour without truncation. Keep the harness: in
\cref{ch:ch13} you will add RVO and \orca to the comparison.
\end{exercise}

\begin{exercise}[\difficulty{2}]
\label{exr:ch12-3d}
(a)~Prove \cref{thm:ch12-cone3d}: the set of vectors $\vel \in \R^3$ whose ray
meets the ball $\disc{\pos_{\mathrm{rel}}}{R}$ is the right circular cone of
half-angle $\arcsin(R/d)$ about $\pos_{\mathrm{rel}}$, and its intersection
with any plane through the axis is the planar cone of \cref{thm:ch12-cone}.
(b)~Show that the cone touches the sphere along a circle and compute its
centre and radius in terms of $d$ and $R$. (c)~Extend \code{velocity\_samples}
to three dimensions (speeds times well-spread directions on the unit sphere)
and, with \code{time\_to\_collision}, which already accepts 3-vectors, check in
a head-on encounter that a drone can escape by climbing, and compare the
deviation from the preferred velocity with the planar escape.
\end{exercise}
